\chapter{CHUYỂN ĐỔI SANG LƯỢC ĐỒ QUAN HỆ VÀ CHUẨN HÓA BCNF}

\section{Quy tắc ánh xạ từ Lược đồ E-R sang Mô hình Quan hệ}

Quá trình chuyển đổi từ mô hình thực thể liên kết (E-R) sang mô hình dữ liệu quan hệ (Relational Model) được thực hiện tuân thủ theo các nguyên tắc chuẩn tắc của lý thuyết cơ sở dữ liệu:
\begin{enumerate}
    \item \textbf{Ánh xạ tập thực thể mạnh}: Mỗi tập thực thể mạnh độc lập (\texttt{SAN\_VAN\_DONG}, \texttt{DOI\_BONG}, \texttt{CAU\_THU}, \texttt{TRONG\_TAI}, \texttt{VONG\_DAU}) chuyển đổi thành một quan hệ độc lập. Khóa chính của thực thể trở thành khóa chính (PK) của quan hệ.
    \item \textbf{Ánh xạ mối kết hợp $1 - N$}: Mối kết hợp một - nhiều được chuyển đổi bằng cách lấy khóa chính của thực thể phía bên một ($1$) đưa vào làm thuộc tính khóa ngoại (FK) tại quan hệ phía bên nhiều ($N$). Ví dụ: Thuộc tính \texttt{MaSanNha} (FK) trong quan hệ \texttt{DOI\_BONG} tham chiếu đến \texttt{MaSan} trong \texttt{SAN\_VAN\_DONG}.
    \item \textbf{Ánh xạ thuộc tính đa trị và thực thể yếu}: Thuộc tính màu áo của đội bóng (mỗi đội có $\ge 2$ màu áo) được ánh xạ thành quan hệ \texttt{MAU\_AO\_DOI\_BONG}. Khóa chính tổ hợp gồm khóa ngoại của thực thể chủ \texttt{MaDoi} và thuộc tính phân biệt \texttt{LoaiMauAo}.
    \item \textbf{Ánh xạ mối kết hợp $N - N$}: Mối kết hợp nhiều - nhiều giữa \texttt{TRAN\_DAU} và \texttt{TRONG\_TAI} được tách thành quan hệ trung gian \texttt{PHAN\_CONG\_TRONG\_TAI} với khóa chính tổ hợp gồm \texttt{MaTran} và \texttt{VaiTro}. Tương tự, danh sách thi đấu được biểu diễn qua quan hệ \texttt{DANG\_KY\_THI\_DAU}.
\end{enumerate}

---

\section{Tập hợp các Lược đồ Quan hệ và Liên kết Khóa}

Sau quá trình ánh xạ, hệ thống bao gồm 11 lược đồ quan hệ. Ký hiệu quy ước: Các thuộc tính \underline{\textbf{in đậm gạch chân}} biểu diễn \textbf{Khóa chính (PK)}, các thuộc tính \textit{in nghiêng} biểu diễn \textbf{Khóa ngoại (FK)}.

\begin{enumerate}
    \item $\text{SAN\_VAN\_DONG}(\underline{\mathbf{MaSan}}, \text{TenSan}, \text{DiaChi}, \text{SucChua})$
    
    \item $\text{DOI\_BONG}(\underline{\mathbf{MaDoi}}, \text{TenDoi}, \text{HLV}, \textit{MaSanNha})$\\
    \textit{Khóa ngoại: $\text{MaSanNha} \rightarrow \text{SAN\_VAN\_DONG}(\text{MaSan})$}
    
    \item $\text{MAU\_AO\_DOI\_BONG}(\underline{\mathbf{MaDoi}}, \underline{\mathbf{LoaiMauAo}}, \text{MoTaMau})$\\
    \textit{Khóa ngoại: $\text{MaDoi} \rightarrow \text{DOI\_BONG}(\text{MaDoi})$}
    
    \item $\text{CAU\_THU}(\underline{\mathbf{MaCauThu}}, \text{HoTen}, \text{NgaySinh}, \text{ViTri}, \text{QuocTich})$
    
    \item $\text{HOP\_DONG\_CAU\_THU}(\underline{\mathbf{MaHD}}, \textit{MaCauThu}, \textit{MaDoiChuQuan}, \textit{MaDoiHienTai}, \text{SoAo}, \text{NgayBatDau}, \text{NgayKetThuc}, \text{TrangThai}, \text{MucLuong})$\\
    \textit{Khóa ngoại: $\text{MaCauThu} \rightarrow \text{CAU\_THU}(\text{MaCauThu})$}\\
    \textit{Khóa ngoại: $\text{MaDoiChuQuan} \rightarrow \text{DOI\_BONG}(\text{MaDoi})$}\\
    \textit{Khóa ngoại: $\text{MaDoiHienTai} \rightarrow \text{DOI\_BONG}(\text{MaDoi})$}
    
    \item $\text{TRONG\_TAI}(\underline{\mathbf{MaTrongTai}}, \text{HoTen}, \text{NgaySinh}, \text{QueQuan}, \text{CapBac})$
    
    \item $\text{VONG\_DAU}(\underline{\mathbf{MaVong}}, \text{TenVong}, \text{ThuTu})$
    
    \item $\text{TRAN\_DAU}(\underline{\mathbf{MaTran}}, \textit{MaVong}, \textit{MaDoiNha}, \textit{MaDoiKhach}, \textit{MaSan}, \text{NgayGio}, \text{GiaVe}, \text{SoKhach}, \text{BanThangNha}, \text{BanThangKhach}, \text{TheVangNha}, \text{TheVangKhach}, \text{TheDoNha}, \text{TheDoKhach}, \text{TrangThai})$\\
    \textit{Khóa ngoại: $\text{MaVong} \rightarrow \text{VONG\_DAU}(\text{MaVong})$}\\
    \textit{Khóa ngoại: $\text{MaDoiNha} \rightarrow \text{DOI\_BONG}(\text{MaDoi})$}\\
    \textit{Khóa ngoại: $\text{MaDoiKhach} \rightarrow \text{DOI\_BONG}(\text{MaDoi})$}\\
    \textit{Khóa ngoại: $\text{MaSan} \rightarrow \text{SAN\_VAN\_DONG}(\text{MaSan})$}
    
    \item $\text{PHAN\_CONG\_TRONG\_TAI}(\underline{\mathbf{MaTran}}, \underline{\mathbf{VaiTro}}, \textit{MaTrongTai})$\\
    \textit{Khóa ngoại: $\text{MaTran} \rightarrow \text{TRAN\_DAU}(\text{MaTran})$}\\
    \textit{Khóa ngoại: $\text{MaTrongTai} \rightarrow \text{TRONG\_TAI}(\text{MaTrongTai})$}
    
    \item $\text{DANG\_KY\_THI\_DAU}(\underline{\mathbf{MaTran}}, \underline{\mathbf{MaCauThu}}, \textit{MaDoi}, \text{LoaiDangKy}, \text{SoAoTran})$\\
    \textit{Khóa ngoại: $\text{MaTran} \rightarrow \text{TRAN\_DAU}(\text{MaTran})$}\\
    \textit{Khóa ngoại: $\text{MaCauThu} \rightarrow \text{CAU\_THU}(\text{MaCauThu})$}\\
    \textit{Khóa ngoại: $\text{MaDoi} \rightarrow \text{DOI\_BONG}(\text{MaDoi})$}
    
    \item $\text{SU\_KIEN\_TRAN\_DAU}(\underline{\mathbf{MaSuKien}}, \textit{MaTran}, \textit{MaCauThu}, \textit{MaDoi}, \text{LoaiSuKien}, \text{PhutGhiNhan})$\\
    \textit{Khóa ngoại: $\text{MaTran} \rightarrow \text{TRAN\_DAU}(\text{MaTran})$}\\
    \textit{Khóa ngoại: $\text{MaCauThu} \rightarrow \text{CAU\_THU}(\text{MaCauThu})$}\\
    \textit{Khóa ngoại: $\text{MaDoi} \rightarrow \text{DOI\_BONG}(\text{MaDoi})$}
\end{enumerate}

\begin{figure}[H]
    \centering
    \fbox{\parbox{0.9\textwidth}{\centering \vspace{3cm} \textbf{[SƠ ĐỒ QUAN HỆ VẬT LÝ VÀ LIÊN KẾT KHÓA CHÍNH - NGOẠI]} \\ \textit{(Tham khảo mã Mermaid tại file docs/SO_DO_MERMAID_DRAWIO.md)} \vspace{3cm}}}
    \caption{Sơ đồ cấu trúc lược đồ quan hệ và mối liên kết khóa}
    \label{fig:relational_diagram}
\end{figure}

---

\section{Cơ sở lý thuyết Phụ thuộc hàm và Dạng chuẩn BCNF}

\subsection{Khái niệm Phụ thuộc hàm}
Cho lược đồ quan hệ $R(U)$ với tập thuộc tính $U$ và $X, Y \subseteq U$. Một \textbf{phụ thuộc hàm (Functional Dependency - FD)}, ký hiệu $X \rightarrow Y$, biểu thị rằng: với mọi bộ giá trị $t_1, t_2$ trong quan hệ $r(R)$, nếu $t_1[X] = t_2[X]$ thì bắt buộc $t_1[Y] = t_2[Y]$.
\begin{itemize}
    \item Phụ thuộc hàm được gọi là \textit{tầm thường} nếu $Y \subseteq X$.
    \item Phụ thuộc hàm được gọi là \textit{không tầm thường} nếu $Y \not\subseteq X$.
\end{itemize}

\subsection{Bao đóng thuộc tính và Khóa ứng viên}
Cho tập phụ thuộc hàm $F$. Bao đóng của tập thuộc tính $X$ dưới tập $F$, ký hiệu $X^+$, là tập hợp tất cả các thuộc tính $A \in U$ sao cho phụ thuộc hàm $X \rightarrow A$ được suy dẫn logic từ $F$ (dựa trên hệ tiên đề Armstrong).
\begin{itemize}
    \item Tập thuộc tính $X$ là một \textbf{siêu khóa (Superkey)} của quan hệ $R(U)$ nếu $X^+ = U$.
    \item Tập thuộc tính $X$ là một \textbf{khóa ứng viên (Candidate Key)} nếu $X$ là siêu khóa và không tồn tại bất kỳ tập con thực sự $X' \subset X$ nào sao cho $X'^+ = U$ (tính tối thiểu).
\end{itemize}

\subsection{Định nghĩa Dạng chuẩn Boyce-Codd (BCNF)}
Một lược đồ quan hệ $R(U)$ cùng với tập phụ thuộc hàm $F$ được gọi là đạt \textbf{Dạng chuẩn Boyce-Codd (BCNF)} nếu và chỉ nếu:
$$\forall (X \rightarrow Y) \in F^+ \text{ với } Y \not\subseteq X \implies X \text{ là một siêu khóa của } R.$$

Dạng chuẩn BCNF là một trong những dạng chuẩn khắt khe nhất, loại bỏ triệt để mọi dư thừa dữ liệu có thể suy diễn từ phụ thuộc hàm. Mọi quan hệ đạt BCNF mặc nhiên đạt dạng chuẩn 3NF, 2NF và 1NF.

---

\section{Xác định Phụ thuộc hàm và Chứng minh Chuẩn BCNF cho từng Lược đồ}

\subsection{Quan hệ SAN\_VAN\_DONG}
\begin{itemize}
    \item Tập thuộc tính: $U = \{\text{MaSan}, \text{TenSan}, \text{DiaChi}, \text{SucChua}\}$.
    \item Tập phụ thuộc hàm tối thiểu $F_1$:
    \begin{align*}
        f_{1.1}: & \quad \text{MaSan} \rightarrow \{\text{TenSan}, \text{DiaChi}, \text{SucChua}\} \\
        f_{1.2}: & \quad \text{TenSan} \rightarrow \{\text{MaSan}, \text{DiaChi}, \text{SucChua}\}
    \end{align*}
    \item Xác định khóa ứng viên:
    $$\text{MaSan}^+ = U \implies \text{MaSan} \text{ là khóa ứng viên.}$$
    $$\text{TenSan}^+ = U \implies \text{TenSan} \text{ là khóa ứng viên.}$$
    \item Kiểm tra BCNF:
    \begin{itemize}
        \item Xét $f_{1.1}$: Vế trái $\text{MaSan}$ là siêu khóa.
        \item Xét $f_{1.2}$: Vế trái $\text{TenSan}$ là siêu khóa.
    \end{itemize}
    \item \textbf{Kết luận}: Quan hệ \texttt{SAN\_VAN\_DONG} đạt dạng chuẩn **BCNF**.
\end{itemize}

\subsection{Quan hệ DOI\_BONG}
\begin{itemize}
    \item Tập thuộc tính: $U = \{\text{MaDoi}, \text{TenDoi}, \text{HLV}, \text{MaSanNha}\}$.
    \item Tập phụ thuộc hàm $F_2$:
    \begin{align*}
        f_{2.1}: & \quad \text{MaDoi} \rightarrow \{\text{TenDoi}, \text{HLV}, \text{MaSanNha}\} \\
        f_{2.2}: & \quad \text{TenDoi} \rightarrow \{\text{MaDoi}, \text{HLV}, \text{MaSanNha}\}
    \end{align*}
    \item Xác định khóa: Khóa ứng viên là $\{\text{MaDoi}\}$ và $\{\text{TenDoi}\}$.
    \item Kiểm tra BCNF: Cả hai vế trái $\text{MaDoi}$ và $\text{TenDoi}$ đều là siêu khóa của quan hệ.
    \item \textbf{Kết luận}: Quan hệ \texttt{DOI\_BONG} đạt dạng chuẩn **BCNF**.
\end{itemize}

\subsection{Quan hệ MAU\_AO\_DOI\_BONG}
\begin{itemize}
    \item Tập thuộc tính: $U = \{\text{MaDoi}, \text{LoaiMauAo}, \text{MoTaMau}\}$.
    \item Tập phụ thuộc hàm $F_3$:
    $$f_{3.1}: \quad \{\text{MaDoi}, \text{LoaiMauAo}\} \rightarrow \text{MoTaMau}$$
    \item Khóa ứng viên: $K = \{\text{MaDoi}, \text{LoaiMauAo}\}$ do $\{\text{MaDoi}, \text{LoaiMauAo}\}^+ = U$.
    \item Kiểm tra BCNF: Vế trái của phụ thuộc hàm duy nhất $f_{3.1}$ chính là khóa chính (siêu khóa).
    \item \textbf{Kết luận}: Quan hệ \texttt{MAU\_AO\_DOI\_BONG} đạt dạng chuẩn **BCNF**.
\end{itemize}

\subsection{Quan hệ CAU\_THU}
\begin{itemize}
    \item Tập thuộc tính: $U = \{\text{MaCauThu}, \text{HoTen}, \text{NgaySinh}, \text{ViTri}, \text{QuocTich}\}$.
    \item Tập phụ thuộc hàm $F_4$:
    $$f_{4.1}: \quad \text{MaCauThu} \rightarrow \{\text{HoTen}, \text{NgaySinh}, \text{ViTri}, \text{QuocTich}\}$$
    \item Khóa ứng viên: $K = \{\text{MaCauThu}\}$.
    \item Kiểm tra BCNF: Vế trái $\text{MaCauThu}$ là siêu khóa duy nhất.
    \item \textbf{Kết luận}: Quan hệ \texttt{CAU\_THU} đạt dạng chuẩn **BCNF**.
\end{itemize}

\subsection{Quan hệ HOP\_DONG\_CAU\_THU}
\begin{itemize}
    \item Tập thuộc tính: $U = \{\text{MaHD}, \text{MaCauThu}, \text{MaDoiChuQuan}, \text{MaDoiHienTai}, \text{SoAo}, \text{NgayBatDau}, \text{NgayKetThuc}, \text{TrangThai}, \text{MucLuong}\}$.
    \item Tập phụ thuộc hàm $F_5$:
    $$f_{5.1}: \quad \text{MaHD} \rightarrow \{\text{MaCauThu}, \text{MaDoiChuQuan}, \text{MaDoiHienTai}, \text{SoAo}, \text{NgayBatDau}, \text{NgayKetThuc}, \text{TrangThai}, \text{MucLuong}\}$$
    \item Khóa ứng viên: $K = \{\text{MaHD}\}$.
    \item Kiểm tra BCNF: Vế trái $\text{MaHD}$ là siêu khóa. Không có phụ thuộc bắc cầu hay phụ thuộc một phần.
    \item \textbf{Kết luận}: Quan hệ \texttt{HOP\_DONG\_CAU\_THU} đạt dạng chuẩn **BCNF**.
\end{itemize}

\subsection{Quan hệ TRONG\_TAI}
\begin{itemize}
    \item Tập thuộc tính: $U = \{\text{MaTrongTai}, \text{HoTen}, \text{NgaySinh}, \text{QueQuan}, \text{CapBac}\}$.
    \item Tập phụ thuộc hàm $F_6$:
    $$f_{6.1}: \quad \text{MaTrongTai} \rightarrow \{\text{HoTen}, \text{NgaySinh}, \text{QueQuan}, \text{CapBac}\}$$
    \item Khóa ứng viên: $K = \{\text{MaTrongTai}\}$.
    \item Kiểm tra BCNF: Vế trái $\text{MaTrongTai}$ là siêu khóa.
    \item \textbf{Kết luận}: Quan hệ \texttt{TRONG\_TAI} đạt dạng chuẩn **BCNF**.
\end{itemize}

\subsection{Quan hệ VONG\_DAU}
\begin{itemize}
    \item Tập thuộc tính: $U = \{\text{MaVong}, \text{TenVong}, \text{ThuTu}\}$.
    \item Tập phụ thuộc hàm $F_7$:
    \begin{align*}
        f_{7.1}: & \quad \text{MaVong} \rightarrow \{\text{TenVong}, \text{ThuTu}\} \\
        f_{7.2}: & \quad \text{ThuTu} \rightarrow \{\text{MaVong}, \text{TenVong}\}
    \end{align*}
    \item Khóa ứng viên: $K_1 = \{\text{MaVong}\}$, $K_2 = \{\text{ThuTu}\}$.
    \item Kiểm tra BCNF: Vế trái của mọi phụ thuộc hàm không tầm thường đều là siêu khóa.
    \item \textbf{Kết luận}: Quan hệ \texttt{VONG\_DAU} đạt dạng chuẩn **BCNF**.
\end{itemize}

\subsection{Quan hệ TRAN\_DAU}
\begin{itemize}
    \item Tập thuộc tính:
    $$U = \{\text{MaTran}, \text{MaVong}, \text{MaDoiNha}, \text{MaDoiKhach}, \text{MaSan}, \text{NgayGio}, \text{GiaVe}, \text{SoKhach}, \text{BanThangNha}, \text{BanThangKhach}, \dots\}$$
    \item Tập phụ thuộc hàm $F_8$:
    \begin{align*}
        f_{8.1}: & \quad \text{MaTran} \rightarrow U \setminus \{\text{MaTran}\} \\
        f_{8.2}: & \quad \{\text{MaDoiNha}, \text{MaDoiKhach}\} \rightarrow U \setminus \{\text{MaDoiNha}, \text{MaDoiKhach}\}
    \end{align*}
    \textit{(Do thể thức vòng tròn 2 lượt quy định mỗi cặp đối đầu sân nhà - sân khách chỉ diễn ra duy nhất 1 trận trong mùa giải)}.
    \item Khóa ứng viên: $K_1 = \{\text{MaTran}\}$, $K_2 = \{\text{MaDoiNha}, \text{MaDoiKhach}\}$.
    \item Kiểm tra BCNF: Cả hai vế trái $\text{MaTran}$ và $\{\text{MaDoiNha}, \text{MaDoiKhach}\}$ đều là siêu khóa.
    \item \textbf{Kết luận}: Quan hệ \texttt{TRAN\_DAU} đạt dạng chuẩn **BCNF**.
\end{itemize}

\subsection{Quan hệ PHAN\_CONG\_TRONG\_TAI}
\begin{itemize}
    \item Tập thuộc tính: $U = \{\text{MaTran}, \text{VaiTro}, \text{MaTrongTai}\}$.
    \item Tập phụ thuộc hàm $F_9$:
    \begin{align*}
        f_{9.1}: & \quad \{\text{MaTran}, \text{VaiTro}\} \rightarrow \text{MaTrongTai} \\
        f_{9.2}: & \quad \{\text{MaTran}, \text{MaTrongTai}\} \rightarrow \text{VaiTro}
    \end{align*}
    \textit{(Một trận đấu với 1 vai trò chỉ có đúng 1 trọng tài; và trong 1 trận, 1 trọng tài chỉ đảm nhiệm tối đa 1 vai trò)}.
    \item Khóa ứng viên: $K_1 = \{\text{MaTran}, \text{VaiTro}\}$, $K_2 = \{\text{MaTran}, \text{MaTrongTai}\}$.
    \item Kiểm tra BCNF: Vế trái của $f_{9.1}$ và $f_{9.2}$ đều là siêu khóa của quan hệ.
    \item \textbf{Kết luận}: Quan hệ \texttt{PHAN\_CONG\_TRONG\_TAI} đạt dạng chuẩn **BCNF**.
\end{itemize}

\subsection{Quan hệ DANG\_KY\_THI\_DAU}
\begin{itemize}
    \item Tập thuộc tính: $U = \{\text{MaTran}, \text{MaCauThu}, \text{MaDoi}, \text{LoaiDangKy}, \text{SoAoTran}\}$.
    \item Tập phụ thuộc hàm $F_{10}$:
    $$f_{10.1}: \quad \{\text{MaTran}, \text{MaCauThu}\} \rightarrow \{\text{MaDoi}, \text{LoaiDangKy}, \text{SoAoTran}\}$$
    \item Khóa ứng viên: $K = \{\text{MaTran}, \text{MaCauThu}\}$.
    \item Kiểm tra BCNF: Vế trái $\{\text{MaTran}, \text{MaCauThu}\}$ là siêu khóa duy nhất.
    \item \textbf{Kết luận}: Quan hệ \texttt{DANG\_KY\_THI\_DAU} đạt dạng chuẩn **BCNF**.
\end{itemize}

\subsection{Quan hệ SU\_KIEN\_TRAN\_DAU}
\begin{itemize}
    \item Tập thuộc tính: $U = \{\text{MaSuKien}, \text{MaTran}, \text{MaCauThu}, \text{MaDoi}, \text{LoaiSuKien}, \text{PhutGhiNhan}\}$.
    \item Tập phụ thuộc hàm $F_{11}$:
    $$f_{11.1}: \quad \text{MaSuKien} \rightarrow \{\text{MaTran}, \text{MaCauThu}, \text{MaDoi}, \text{LoaiSuKien}, \text{PhutGhiNhan}\}$$
    \item Khóa ứng viên: $K = \{\text{MaSuKien}\}$.
    \item Kiểm tra BCNF: Vế trái $\text{MaSuKien}$ là khóa chính (siêu khóa).
    \item \textbf{Kết luận}: Quan hệ \texttt{SU\_KIEN\_TRAN\_DAU} đạt dạng chuẩn **BCNF**.
\end{itemize}

---

\section{Phân tích Thuật toán Tách chuẩn BCNF (Phản ví dụ thực tế)}

Để chứng minh tính chuẩn xác và chiều sâu học thuật của mô hình thiết kế, mục này trình bày một ví dụ thiết kế chưa chuẩn hóa và áp dụng từng bước của **Thuật toán tách chuẩn BCNF (BCNF Decomposition Algorithm)**.

\subsection{Tình huống thiết kế vi phạm chuẩn hóa}
Giả sử người thiết kế gộp chung thông tin của Trận đấu và Sân vận động vào trong một bảng duy nhất:
$$R_{\text{Gop}}(\underline{\text{MaTran}}, \text{NgayGio}, \text{MaSan}, \text{TenSan}, \text{DiaChi}, \text{SucChua})$$

Tập phụ thuộc hàm trên $R_{\text{Gop}}$ gồm có:
\begin{align*}
    f_1: & \quad \text{MaTran} \rightarrow \{\text{NgayGio}, \text{MaSan}, \text{TenSan}, \text{DiaChi}, \text{SucChua}\} \\
    f_2: & \quad \text{MaSan} \rightarrow \{\text{TenSan}, \text{DiaChi}, \text{SucChua}\}
\end{align*}

\paragraph{Phân tích vi phạm BCNF:}
\begin{itemize}
    \item Khóa chính duy nhất của $R_{\text{Gop}}$ là $\{\text{MaTran}\}$.
    \item Xét phụ thuộc hàm $f_2: \text{MaSan} \rightarrow \{\text{TenSan}, \text{DiaChi}, \text{SucChua}\}$. Vế trái $\text{MaSan}$ \textbf{không phải là siêu khóa} của $R_{\text{Gop}}$ (vì $\text{MaSan}^+ = \{\text{MaSan}, \text{TenSan}, \text{DiaChi}, \text{SucChua}\} \neq U$).
    \item Do đó, quan hệ $R_{\text{Gop}}$ vi phạm nghiêm trọng định nghĩa BCNF (và cũng vi phạm 3NF do phụ thuộc bắc cầu).
    \item \textbf{Hậu quả thực tế}: Thông tin tên sân, địa chỉ, sức chứa bị lặp đi lặp lại ở mọi trận đấu diễn ra trên sân đó; khi sân vận động thay đổi sức chứa thì phải cập nhật nhiều dòng; nếu chưa có trận đấu nào trên sân thì không thể lưu thông tin sân (Dị thường thêm, sửa, xóa).
\end{itemize}

\subsection{Các bước thực thi thuật toán tách BCNF}
Thuật toán tách chuẩn BCNF tiến hành như sau:
\begin{itemize}
    \item \textbf{Bước 1}: Tìm một phụ thuộc hàm $X \rightarrow Y$ vi phạm điều kiện BCNF. Ta chọn $f_2: \text{MaSan} \rightarrow \{\text{TenSan}, \text{DiaChi}, \text{SucChua}\}$.
    \item \textbf{Bước 2}: Phân rã $R_{\text{Gop}}$ thành hai lược đồ quan hệ con $R_1$ và $R_2$:
    $$R_1 = X \cup Y = \{\underline{\mathbf{MaSan}}, \text{TenSan}, \text{DiaChi}, \text{SucChua}\} \equiv \text{SAN\_VAN\_DONG}$$
    $$R_2 = R_{\text{Gop}} \setminus (Y \setminus X) = \{\underline{\mathbf{MaTran}}, \text{NgayGio}, \textit{MaSan}\} \equiv \text{TRAN\_DAU\_RUT\_GON}$$
    \item \textbf{Bước 3}: Kiểm tra tính chất của phép phân rã:
    \begin{enumerate}
        \item \textit{Tính kết nối không mất mát thông tin (Lossless Join)}:
        $$R_1 \cap R_2 = \{\text{MaSan}\}$$
        Vì $\text{MaSan} \rightarrow R_1$ nên giao $R_1 \cap R_2$ là siêu khóa của $R_1$. Theo định lý kết nối không mất mát, phép phân rã bảo toàn $100\%$ dữ liệu ban đầu:
        $$R_{\text{Gop}} = R_1 \bowtie R_2.$$
        \item \textit{Tính bảo toàn phụ thuộc hàm}: Mọi phụ thuộc hàm trong tập $F$ ban đầu đều được duy trì trong các bảng phân rã con.
        \item \textit{Kiểm tra lại BCNF}: Cả $R_1$ và $R_2$ đều thỏa mãn điều kiện BCNF.
    \end{enumerate}
\end{itemize}

Kết quả chứng minh rằng toàn bộ 11 bảng của hệ thống được thiết kế hoàn toàn tối ưu, loại bỏ triệt để mọi nguy cơ dị thường dữ liệu.
